\documentclass[11pt]{article}

\usepackage{amsmath,amssymb,amsthm}
\usepackage[margin=1in]{geometry}
\usepackage{xurl}
\usepackage{hyperref}
\hypersetup{
  pdftitle={A Hahn-Series Coefficientwise Proof of the Hickerson--Mortenson Hecke--Appell Identities},
  pdfauthor={Xiang Huang}
}

\newtheorem{theorem}{Theorem}[section]
\newtheorem{proposition}[theorem]{Proposition}
\newtheorem{corollary}[theorem]{Corollary}
\newtheorem{lemma}[theorem]{Lemma}
\newtheorem{remark}[theorem]{Remark}
\newtheorem{definition}[theorem]{Definition}

\newcommand{\Z}{\mathbb{Z}}
\newcommand{\Q}{\mathbb{Q}}
\newcommand{\sg}{\operatorname{sg}}
\newcommand{\sgtwo}{\operatorname{sg}_2}

\title{A Hahn-Series Coefficientwise Proof of the\\
Hickerson--Mortenson Hecke--Appell Identities}
\author{Xiang Huang}
\date{August 2026}

\begin{document}

\maketitle

\begin{abstract}
We give a coefficientwise proof guided by an ordered Hahn-series expansion
of the Hickerson--Mortenson identity expressing the Hecke-type
double sum $f_{n,n+p,n}$ as an Appell--Lerch sum plus a theta
correction, for all $n,p\ge 1$.  The decomposition itself is derived without convergence estimates or
Ramanujan's ${}_1\psi_1$ summation, using a single algebraic
mechanism: the Hahn order assigns each Appell--Lerch denominator
$1/(1-M)$ a branch sign $\sg(\gamma)$ determined by
the monomial order, and the resulting coefficient-level sign-wall kernel $\sgtwo$
is the theta correction---not an ad hoc residue, but the forced
cost of the chosen expansion chamber.  A row-exponent ring
identity collapses the multi-index Appell exponent to the Hecke
quadratic $H(r,s)$ plus a pure theta shell $N\binom{w+1}{2}$.
The remaining wall term is identified with the
Hickerson--Mortenson theta correction $\theta_{n,p}$ via a
coefficientwise Hahn proof of the Kronecker partial-fraction identity.
At the paper level, the identity is a family of $\Z$-coefficient
identities, one per monomial.  The Lean statements realize these
coefficient kernels over $\Q$; integrality and general-ring base change are
not separate formalized theorems.  The general coefficient theorem, the
source-faithful comparison with HM's printed correction, and the
coefficientwise Kronecker argument have been verified in Lean~4 with zero
\texttt{sorry}
and no axioms beyond \texttt{propext},
\texttt{Classical.choice}, and \texttt{Quot.sound}.
\end{abstract}

%----------------------------------------------------------------------
\section{Introduction}
%----------------------------------------------------------------------

\subsection*{Motivation: bilateral algebra needs an order}

Naive bilateral formal manipulation is inconsistent.
In a Laurent ring, the geometric identity
$(1-z)\sum_{k\ge 0}z^k=1$ is unambiguous, but a bilateral
``identity'' $(1-z)\sum_{k\in\Z}z^k=0$ would yield $1=0$ by
dividing.  The root cause is that $\sum_{k\in\Z}z^k$ has no
well-ordered support.  In a Hahn-series ring with a
lexicographic monomial order, the inverse $(1-M)^{-1}$ is
uniquely defined: it expands as $\sum_{\ell\ge 0}M^\ell$
or $-\sum_{\ell\ge 1}M^{-\ell}$, depending on the sign of the
leading exponent of~$M$.  The branch choice is forced, and
the branch sign is the price the algebra pays for
well-definedness.

This paper shows that the Hickerson--Mortenson theta correction
in a Hecke-type identity is exactly that price: the sign-wall
kernel $\sgtwo$ produced by expanding Appell--Lerch denominators
in the Hahn ring.

\subsection*{Background}

Hickerson and Mortenson~\cite{jHiMo14} proved that every Hecke-type
double sum $f_{n,n+p,n}$ (stated under $\gcd(n,p)=1$) decomposes as an
Appell--Lerch sum $g_{n,n+p,n}$ plus a theta correction
$\theta_{n,p}/\bar{J}_{0,N}$, where $N=np(2n+p)$.  Their proof uses
analytic functional equations and a growth-based uniqueness argument.
Mortenson~\cite{jMo12} gave an alternative proof using Ramanujan's
${}_1\psi_1$ summation.  Both proofs rely on convergence in the region
$|q|<1$ and produce the theta correction as an analytic residue.
For background on mock theta functions and
Appell--Lerch sums, see Zwegers~\cite{oZw02} and Chan~\cite{oCh11};
for Bailey-pair approaches to Hecke-type double sums producing
one-variable specializations, see Andrews~\cite{jAn01} and
Lovejoy--Osburn~\cite{jLO13}.

\subsection*{This paper}

We give a Hahn-series coefficientwise proof of the decomposition
that replaces the analytic ingredients by formal algebra.  The key observation
is that the Hahn order assigns each Appell--Lerch denominator
$1/(1-M)$ a branch sign $\sg(\gamma)$ determined by the
monomial order.  After clearing the $\bar{J}$-factor and expanding all
Appell rows, the resulting coefficient is a sign-wall function
$\sgtwo$ that differs from the Hecke cone indicator by exactly
the theta correction $\bar\Theta_{n,p}$.

The proof reduces to two ingredients:
\begin{enumerate}
\item A \emph{row-exponent ring identity} (Proposition~\ref{prop:exponent})
  showing that the Appell row exponent simplifies to $H(r,s)+N\binom{w+1}{2}$.
\item The \emph{signed Hahn geometric inverse} (Lemma~\ref{lem:hahn})
  producing the branch coefficient $(\sg(s)-\sg(V_w))/2$ per row.
\end{enumerate}
The displayed kernels are integer-valued, so the paper argument gives
$\Z$-identities one coefficient at a time.  The Lean theorem is stated over
$\Q$ and does not separately package integrality or base change to an
arbitrary commutative ring.  The coefficient data is nevertheless available
as input to radial-limit and quantum-modular analysis~\cite{jZag10,jFOR13},
where $|q|<1$ proofs do not apply directly.

Table~\ref{tab:replacement} summarizes how each analytic ingredient
is replaced.

\begin{table}[ht]
\centering
\begin{tabular}{ll}
\hline
\textbf{Analytic ingredient} & \textbf{Hahn-series replacement} \\
\hline
Convergence ($|q|<|x|<1$) & Lexicographic monomial order \\
Geometric expansion branch & Signed Hahn inverse (Lemma~\ref{lem:hahn}) \\
Analytic continuation & Coefficientwise identity \\
Poles, residues, uniqueness & Coefficientwise kernel identity \\
Theta correction formula & Sign-wall boundary term ($\sgtwo$) \\
${}_1\psi_1$ summation & Not used (Appendix uses Kronecker PF) \\
\hline
\end{tabular}
\caption{Analytic ingredients and their formal replacements.}
\label{tab:replacement}
\end{table}

\medskip
\noindent\textbf{Scope.}
The coefficientwise proof of the sign-wall identity is
contained in Sections~\ref{sec:defs}--\ref{sec:proof}.
The analytic specialization (recovering HM's analytic identity)
requires matching the Hahn branch with the convergent branch
on the unit torus (Remark~\ref{rem:analytic}), plus analytic
continuation in $x,y$.
The identification with HM's printed theta correction is proved
separately in Appendix~\ref{app:HM-coeff-comparison}.  Its Kronecker step is
proved coefficientwise from the signed Hahn inverse and a reflected boundary
cancellation.  The Lean~4 development verifies this appendix and the full
general-$(n,p)$ coefficient assembly; the earlier $f_{2,3,2}$ proof remains
as an independent regression instance in HM's spiral parametrization.

\medskip
\noindent\textbf{Outline.}
Section~\ref{sec:defs} gives the definitions and sign conventions.
Section~\ref{sec:hahn} proves the Hahn inverse lemma.
Section~\ref{sec:rows} proves the row-coefficient propositions.
Section~\ref{sec:proof} proves the main theorem.
Section~\ref{sec:examples} gives three examples.
Section~\ref{sec:lean} describes the Lean~4 formalization.

%----------------------------------------------------------------------
\section{Definitions}\label{sec:defs}
%----------------------------------------------------------------------

We use the Jacobi theta function
$j(z;q)=\sum_{k\in\Z}(-1)^k q^{\binom{k}{2}}z^k$,
the Appell--Lerch sum
$m(x,q,z)=\frac{1}{j(z;q)}\sum_{k\in\Z}
\frac{(-1)^k q^{\binom{k}{2}}z^k}{1-q^{k-1}xz}$,
the notation $\bar{J}_{0,N}=j(-1;\,q^N)$, and the Euler product
$J_m=\prod_{i\ge 1}(1-q^{mi})$.

\begin{definition}[Hecke-type double sum]
\[
  f_{A,B,C}(x,y;q)
  =\Bigl(\sum_{r,s\ge 0}-\sum_{r,s<0}\Bigr)
    (-1)^{r+s}x^r y^s\,
    q^{A\binom{r}{2}+Brs+C\binom{s}{2}}.
\]
For $f_{n,n+p,n}$: set $B=n+p$,
$\Delta=p(2n+p)$, $N=n\Delta$, and
\[
  H(r,s)=n\tbinom{r}{2}+(n+p)rs+n\tbinom{s}{2}.
\]
Define the sign-wall functions
$\sg(u)=1$ if $u\ge 0$, $\sg(u)=-1$ if $u<0$, and
$\sgtwo(u,v)=(\sg(u)+\sg(v))/2$.  Thus $\sgtwo$ takes values
in $\{-1,0,1\}$: it equals $1$ when both arguments are
nonnegative, $-1$ when both are negative, and $0$ when they
disagree in sign.

\medskip
\begin{center}
\begin{tabular}{c|cc}
$u\backslash v$ & $v\ge 0$ & $v<0$ \\
\hline
$u\ge 0$ & $\sgtwo=1$ & $\sgtwo=0$ \\
$u<0$ & $\sgtwo=0$ & $\sgtwo=-1$
\end{tabular}
\end{center}
\end{definition}

\begin{definition}[Appell--Lerch row decomposition]
Set $A_0=\frac{np(2n+p+1)}{2}$.
For $t=0,\ldots,n-1$, define the $s$-split rows
\[
G_t^{(s)}=(-y)^t\,q^{n\binom{t}{2}}\,
  j(q^{Bt}x;\,q^n)\cdot
  m\!\Bigl(-q^{A_0-t\Delta}\frac{(-y)^n}{(-x)^B},\;q^N,\;-1\Bigr),
\]
and the $r$-split rows
\[
G_t^{(r)}=(-x)^t\,q^{n\binom{t}{2}}\,
  j(q^{Bt}y;\,q^n)\cdot
  m\!\Bigl(-q^{A_0-t\Delta}\frac{(-x)^n}{(-y)^B},\;q^N,\;-1\Bigr).
\]  Define
$g_{n,n+p,n}=\sum_{t=0}^{n-1}(G_t^{(s)}+G_t^{(r)})$.

These parameters match \cite{jHiMo14}, Theorem~0.3,
specialized to $(z_0,z_1)=(-1,-1)$.
\end{definition}

\begin{definition}[Bar-$J$-cleared theta correction]
\[
  \bar\Theta_{n,p}=\sum_{r,s,w\in\Z}
    (-x)^r(-y)^s\,
    \sgtwo\!\Bigl(r+nw+\tfrac{n-1}{2},\;
                   s-nw-\tfrac{n+1}{2}\Bigr)\,
    q^{H(r,s)+N\binom{w+1}{2}}.
\]
\end{definition}

\begin{proposition}[Identification with HM]\label{prop:HM}
With the conventions above, $\bar\Theta_{n,p}$ and
$\theta_{n,p}^{\mathrm{HM}}$ have the same coefficient function for all
$n,p\ge 1$.  Here $\theta_{n,p}^{\mathrm{HM}}$ is the theta correction
in~\cite{jHiMo14}, Theorem~0.3 (stated there under
$\gcd(n,p)=1$), specialized to $(z_0,z_1)=(-1,-1)$,
taking HM's printed formula as the definition of
$\theta_{n,p}^{\mathrm{HM}}$ in the non-coprime case.
\end{proposition}

\begin{proof}
Appendix~\ref{app:HM-coeff-comparison} gives the direct
coefficient comparison.  It uses the coefficientwise Kronecker
partial-fraction identity (Lemma~\ref{lem:kronecker}) to
expand the HM finite theta quotient, then reindexes from the
residue pair $(r^*,s^*)$ to the integer triple $(R,S,w)$.
The exponent collapses to $H(R,S)+N\binom{w+1}{2}$,
and the branch signs produce $\sgtwo(U_w,V_w)$.

\emph{Note:} This identification is logically independent of
the main proof (Theorem~\ref{thm:main}).  It serves to match
$\bar\Theta_{n,p}$ against HM's printed formula; the
Hecke--Appell decomposition itself (Theorem~\ref{thm:main})
is proved without it.
\end{proof}

\begin{remark}[Consistency check]
As a consistency check (using HM's theorem, not part of our
independent proof), compare coefficients in the cleared HM identity
$\bar{J}\cdot f=\bar{J}\cdot g+\theta_{n,p}^{\mathrm{HM}}$
and Theorem~\ref{thm:main}.  Subtracting shows that
$\bar\Theta_{n,p}$ and $\theta_{n,p}^{\mathrm{HM}}$ have the same
coefficient function.
The self-contained proof of this identification is in
Appendix~\ref{app:HM-coeff-comparison}.
\end{remark}

We record the HM theta correction for completeness.  With
$r=r^*+\{(n-1)/2\}$, $s=s^*+\{(n-1)/2\}$:
\begin{align*}
\theta_{n,p}^{\mathrm{HM}}&=\sum_{r^*=0}^{p-1}\sum_{s^*=0}^{p-1}
 q^{n\binom{r-(n-1)/2}{2}+B(r-(n-1)/2)(s+(n+1)/2)+n\binom{s+(n+1)/2}{2}}\\
&\quad\times(-x)^{r-(n-1)/2}(-y)^{s+(n+1)/2}\\
&\quad\times
\frac{J_{p^2(2n+p)}^3\,
  j(-q^{np(s-r)}x^n/y^n;\,q^{np^2})\,
  j(q^{p(2n+p)(r+s)+p(n+p)}x^py^p;\,q^{p^2(2n+p)})}
  {j(q^{p(2n+p)r+p(n+p)/2}(-y)^B/(-x)^n,\,
     q^{p(2n+p)s+p(n+p)/2}(-x)^B/(-y)^n;\,q^{p^2(2n+p)})}.
\end{align*}
Here $j(a,b;Q):=j(a;Q)\,j(b;Q)$.

\begin{theorem}[Main coefficient identity]\label{thm:main}
Let $n,p\ge 1$.  For every $(e,r,s)\in\Z^3$,
\[
  [q^e x^r y^s]\bigl(\bar{J}_{0,N}\cdot f_{n,n+p,n}\bigr)
  =
  [q^e x^r y^s]\bigl(\bar{J}_{0,N}\cdot g_{n,n+p,n}
    +\bar\Theta_{n,p}\bigr).
\]
\end{theorem}

\begin{remark}
Neither Theorem~\ref{thm:main} nor
Proposition~\ref{prop:HM} requires $\gcd(n,p)=1$: the proof
uses only the division algorithm $s=n\ell+t$ (valid for all
$n\ge 1$), polynomial exponent algebra, and the
Kronecker partial-fraction identity, none of which involve
coprimality.  Hickerson--Mortenson~\cite{jHiMo14} state their
Theorem~0.3 under $\gcd(n,p)=1$; our results extend the identity
to all $n,p\ge 1$ by taking HM's printed formula as the
definition of $\theta_{n,p}^{\mathrm{HM}}$ in the non-coprime
case.
\end{remark}

%----------------------------------------------------------------------
\section{The Signed Hahn Geometric Inverse}\label{sec:hahn}
%----------------------------------------------------------------------

\subsection{Coefficient ring}
We work in the Hahn-series ring~\cite{oHa07}
$S=\mathrm{HahnSeries}\,(\Gamma,\Q)$
where $\Gamma=\Z^3$ with the lexicographic order $(e,a,b)$, so that
$q=\mathfrak{q}^{(1,0,0)}$, $x=\mathfrak{q}^{(0,1,0)}$,
$y=\mathfrak{q}^{(0,0,1)}$.  A monomial $q^e x^a y^b$ has leading
exponent $(e,a,b)$, and the sign $\sg(\gamma)$ refers to the
lexicographic sign of $\gamma$.  In particular, for the Appell
denominator monomial $M_{k,t}=q^D\cdot(-y)^n/(-x)^B$, the leading
exponent is $(D,-B,n)$, which is positive iff $D>0$
(since $D$ is the first component).

All identities below are interpreted coefficientwise: for each
fixed $(e,r,s)$, the coefficient $[q^e x^r y^s]$ of every displayed
series is a finite sum.  The Hahn order is used only to define
the branch of each geometric inverse (Lemma~\ref{lem:hahn}).
Lemma~\ref{lem:finite} below ensures that the theta-shell sums
are locally finite.  More precisely, for each displayed bilateral
sum $\sum_k q^{N\binom{k}{2}}/(1-M_{k,t})$, the monomial
forcing (Step~1 in Proposition~\ref{prop:s-row}) pins
$(t,\ell,u)$ from $(r,s)$, and $e$ determines $k$ to at most
two values, so only finitely many terms contribute to each
coefficient.


\begin{remark}[Coefficientwise reading]\label{rem:coeff}
Theorem~\ref{thm:main} is an identity of coefficient functions: for each
fixed $(e,r,s)\in\Z^3$, both displayed coefficients are finite sums and
agree.  The Hahn order determines the branch of each individual geometric
inverse (Lemma~\ref{lem:hahn}); Lemma~\ref{lem:finite} and
Proposition~\ref{prop:s-row} establish the local finiteness needed for the
coefficient calculation.  We do not infer global well-ordered support from
this local statement, and this paper does not supply a Hahn-series wrapper
for the assembled three-variable families.  In particular, the occurrence
of negative $q$-powers by itself would not decide Hahn-series membership.
\end{remark}

\begin{remark}[Analytic specialization]\label{rem:analytic}
The formal identity specializes to HM's analytic identity
in any chamber where $|M_{k,t}|<1$ iff $\gamma_{k,t}>0$.
Since the denominator monomial $M_{k,t}$ has leading $q$-exponent
$\gamma_{k,t}$ and auxiliary-variable factor $(-y)^n/(-x)^B$,
this condition holds on the unit torus $|x|=|y|=1$
(avoiding Appell poles): there $|M_{k,t}|=|q|^{\gamma_{k,t}}$,
so the Hahn branch coincides with the convergent branch and
the identity specializes by absolute convergence.  For
generic $|x|,|y|\neq 1$, the leading $q$-exponent may not
dominate $|M_{k,t}|$, and analytic continuation in $x,y$
is required to recover HM's identity over all of
$\{|q|<1\}\times(\mathbb{C}^*)^2$.
\end{remark}

\begin{lemma}[Hahn geometric inverse]\label{lem:hahn}
Let $\Gamma$ be a linearly ordered abelian group and let
$S=\mathrm{HahnSeries}(\Gamma,R)$ be the ring of Hahn series
with well-ordered support.  For a unit
monomial $M=c\,\mathfrak{q}^\gamma\in S$ with $c\in R^\times$
and $\gamma\neq 0$:
\[
  (1-M)^{-1}=
  \begin{cases}
    \displaystyle\sum_{\ell\ge 0}M^\ell, & \gamma>0,\\[6pt]
    \displaystyle-\sum_{\ell\ge 1}M^{-\ell}, & \gamma<0.
  \end{cases}
\]
The coefficient of $M^\ell$ is therefore $1$ for $\ell\ge 0$ when
$\gamma>0$, and $-1$ for $\ell\le -1$ when $\gamma<0$;
zero otherwise.
\end{lemma}

\begin{proof}
For $\gamma>0$, the support $\{0,\gamma,2\gamma,\ldots\}$ is
well-ordered in~$\Gamma$, so the series lies in~$S$.
For $\gamma<0$, factor $(1-M)^{-1}=-M^{-1}(1-M^{-1})^{-1}$;
since $-\gamma>0$, the factor $(1-M^{-1})^{-1}$ has well-ordered
support, and the product lies in~$S$.
\end{proof}

\begin{lemma}[Branch coefficient]\label{lem:branch}
For a unit monomial $M$ with leading exponent $\gamma\neq 0$ in the
linearly ordered group, the coefficient of $M^\ell$ in $(1-M)^{-1}$ is
\[
  [M^\ell](1-M)^{-1}=\frac{\sg(\ell)+\sg(\gamma)}{2},
\]
where we use the convention $\sg(0)=1$.
\end{lemma}

\begin{proof}
By Lemma~\ref{lem:hahn}: if $\gamma>0$, the coefficient is $1$ for
$\ell\ge 0$ and $0$ for $\ell<0$, which equals $(1+\sg(\ell))/2
=(\sg(\ell)+\sg(\gamma))/2$.  If $\gamma<0$, the coefficient is
$-1$ for $\ell\le -1$ and $0$ for $\ell\ge 0$, which equals
$(-1+\sg(\ell))/2=(\sg(\ell)+\sg(\gamma))/2$.
\end{proof}

\begin{lemma}[Cleared Appell numerator]\label{lem:cleared}
$\bar{J}_{0,N}\,m(a,q^N,-1)
=\sum_{k\in\Z} q^{N\binom{k}{2}}/(1+q^{N(k-1)}a)$.
\end{lemma}

\begin{proof}
With $z=-1$, the denominator of $m(a,q^N,-1)$ is
$j(-1;\,q^N)=\bar{J}_{0,N}$, so the factor cancels.  In the
numerator, $(-1)^k\cdot(-1)^k=1$, and $1-q^{N(k-1)}a\cdot(-1)
=1+q^{N(k-1)}a$, giving the stated bilateral sum.
\end{proof}

%----------------------------------------------------------------------
\section{Row Coefficients}\label{sec:rows}
%----------------------------------------------------------------------

\begin{lemma}[Local finiteness]\label{lem:finite}
For fixed $(e,r,s)$, the equation $e=H(r,s)+N\binom{w+1}{2}$ has
at most two solutions $w\in\Z$ (namely $w$ and $-w-1$ when
$e-H(r,s)=N\binom{w+1}{2}$).  Hence every coefficient extraction
in the proof is a finite sum.
\end{lemma}

\begin{proof}
$\binom{w+1}{2}=(e-H(r,s))/N$ determines $w(w+1)/2$, which has
at most two integer solutions.
\end{proof}

\begin{lemma}[Branch sign]\label{lem:branch-sign}
Let $0\le t<n$, $k\in\Z$, and set
\[
  W_{k,t}=t+\tfrac{n-1}{2}-nk,\qquad
  \gamma_{k,t}=N(k-1)+A_0-t\Delta.
\]
Then $\gamma_{k,t}\neq 0$ and $\sg(\gamma_{k,t})=-\sg(W_{k,t})$.
\end{lemma}

\begin{proof}
Write $A=nk-t\in\Z$, so $W_{k,t}=(n-1)/2-A$ and
$\gamma_{k,t}=\Delta(A-n/2)+np/2$.

If $W_{k,t}\ge 0$: then $A\le(n-1)/2$, so
$\gamma_{k,t}\le-p(n+p)/2<0$, so $\sg(\gamma_{k,t})=-1=-\sg(W_{k,t})$.

If $W_{k,t}<0$: then $A>(n-1)/2$.  Since $A\in\Z$:
for $n=2m$, $A\ge m$, so $\gamma_{k,t}\ge np/2>0$;
for $n=2m+1$, $A\ge m+1$, so $\gamma_{k,t}\ge\Delta/2+np/2>0$.
In both cases $\sg(\gamma_{k,t})=1=-\sg(W_{k,t})$.
\end{proof}

\begin{proposition}[Row exponent identity]\label{prop:exponent}
For $s=n\ell+t$ with $0\le t<n$:
\[
  n\tbinom{t}{2}+n\tbinom{r+B\ell}{2}+Bt(r+B\ell)
  +N\tbinom{k}{2}+\ell\bigl(N(k-1)+A_0-t\Delta\bigr)
  =H(r,s)+N\tbinom{k+\ell}{2}.
\]
\end{proposition}

\begin{proof}
Both sides are polynomials in $(r,t,\ell,k,n,p)$ of degree~3.
Substituting $s=n\ell+t$, $B=n+p$, $\Delta=B^2-n^2=p(2n+p)$,
$N=n\Delta$, $A_0=np(2n+p+1)/2$, and expanding both sides, the
difference vanishes identically.\footnote{A Sage script
verifying this polynomial identity and the appendix exponent
identity~(A.15) is included as \texttt{check\_exponents.sage}.}
The same identity holds with $r$ and $s$ interchanged.
\end{proof}

\begin{proposition}[$s$-row coefficient]\label{prop:s-row}
Let $V_w=s-nw-(n+1)/2$.  Then
\[
  [q^e\,x^r\,y^s]\;\bar{J}_{0,N}\sum_{t=0}^{n-1}G_t^{(s)}
  =(-1)^{r+s}\sum_{w\in\Z}
    \frac{\sg(s)-\sg(V_w)}{2}\,
    \mathbf{1}_{e=H(r,s)+N\binom{w+1}{2}}.
\]
\end{proposition}

\begin{proof}
Write $s=n\ell+t$ with $0\le t<n$.  The Appell parameter is
$a_t=-q^{A_0-t\Delta}(-y)^n/(-x)^B$.  By Lemma~\ref{lem:cleared},
\begin{equation}\label{eq:cleared-s}
  \bar{J}_{0,N}G_t^{(s)}
  =(-y)^t q^{n\binom{t}{2}}
   \sum_{u\in\Z}(-1)^u q^{n\binom{u}{2}+Btu}\,x^u
   \sum_{k\in\Z}
   \frac{q^{N\binom{k}{2}}}{1-M_{k,t}},
\end{equation}
where
\[
  M_{k,t}=q^{N(k-1)+A_0-t\Delta}\,\frac{(-y)^n}{(-x)^B}
\]
is the denominator monomial (note: $1+q^{N(k-1)}a_t=1-M_{k,t}$).

\emph{Step~1: monomial forcing.}
The monomial $x^r y^s$ from $x^u\cdot M_{k,t}^\ell$ forces
$u-B\ell=r$ (so $u=r+B\ell$) and $n\ell+t=s$ (so $\ell=(s-t)/n$).
Since $M_{k,t}^\ell=q^{\ell\gamma_{k,t}}(-y)^{n\ell}(-x)^{-B\ell}$,
the total sign is $(-1)^t\cdot(-1)^u\cdot(-1)^{n\ell}\cdot(-1)^{-B\ell}
=(-1)^{t+u+n\ell-B\ell}=(-1)^{t+(r+B\ell)+n\ell-B\ell}=(-1)^{r+t+n\ell}
=(-1)^{r+s}$ (using $t+n\ell=s$).

\emph{Step~2: exponent.}
By Proposition~\ref{prop:exponent}, the total $q$-exponent is
$H(r,s)+N\binom{k+\ell}{2}$.  Set $w=k+\ell-1$; then
$\binom{k+\ell}{2}=\binom{w+1}{2}$.

\emph{Step~3: branch coefficient.}
The leading exponent of $M_{k,t}$ in the lexicographic Hahn order
(with $q$ first) is $\gamma_{k,t}=N(k-1)+A_0-t\Delta$.  By
Lemma~\ref{lem:branch}, the coefficient of $M_{k,t}^\ell$ in
$(1-M_{k,t})^{-1}$ is $(\sg(\ell)+\sg(\gamma_{k,t}))/2$.

After substitution $k=w-\ell+1$:
\[
  \gamma_{k,t}=N(w-\ell)+A_0-t\Delta.
\]
By Lemma~\ref{lem:branch-sign} (with $W_{k,t}=V_w$),
$\sg(\gamma_{k,t})=-\sg(V_w)$.  Therefore the branch coefficient is
$(\sg(\ell)-\sg(V_w))/2$.
Since $s=n\ell+t$ with $0\le t<n$ and $\sg(0)=1$, one has
$\sg(\ell)=\sg(s)$, giving the stated formula.
\end{proof}

\begin{proposition}[$r$-row coefficient]\label{prop:r-row}
Let $U_w=r+nw+(n-1)/2$.  Then
\[
  [q^e\,x^r\,y^s]\;\bar{J}_{0,N}\sum_{t=0}^{n-1}G_t^{(r)}
  =(-1)^{r+s}\sum_{w\in\Z}
    \frac{\sg(r)-\sg(U_w)}{2}\,
    \mathbf{1}_{e=H(r,s)+N\binom{w+1}{2}}.
\]
\end{proposition}

\begin{proof}
Write $r=n\ell+t$ with $0\le t<n$.  The Appell parameter is
$a_t=-q^{A_0-t\Delta}(-x)^n/(-y)^B$.  By Lemma~\ref{lem:cleared},
\[
  \bar{J}_{0,N}G_t^{(r)}
  =(-x)^t q^{n\binom{t}{2}}
   \sum_{v\in\Z}(-1)^v q^{n\binom{v}{2}+Btv}\,y^v
   \sum_{k\in\Z}
   \frac{q^{N\binom{k}{2}}}{1-M'_{k,t}},
\]
where $M'_{k,t}=q^{N(k-1)+A_0-t\Delta}(-x)^n/(-y)^B$.

\emph{Monomial forcing:} $y^v\cdot(M'_{k,t})^\ell$ gives $v-B\ell=s$
and $n\ell+t=r$.  Total sign:
$(-1)^{t+v+n\ell-B\ell}=(-1)^{s+t+n\ell}=(-1)^{r+s}$.

\emph{Exponent:} by the symmetric row-exponent identity
(Proposition~\ref{prop:exponent} with $r\leftrightarrow s$),
the $q$-exponent is $H(r,s)+N\binom{k+\ell}{2}$.  Set $w=-k-\ell$
(the opposite sign convention from the $s$-row, chosen so that
$U_w$ matches $W_{k,t}$);
then $\binom{k+\ell}{2}=\binom{-w}{2}=\binom{w+1}{2}$ and
$U_w=r+nw+(n-1)/2=t+(n-1)/2-nk$.

\emph{Branch sign:} by Lemma~\ref{lem:branch-sign} with $W_{k,t}=U_w$,
$\sg(\gamma'_{k,t})=-\sg(U_w)$.  The branch coefficient
is $(\sg(\ell)-\sg(U_w))/2=(\sg(r)-\sg(U_w))/2$.
\end{proof}

\begin{corollary}[Cleared Appell coefficient]\label{cor:g-coeff}
\[
  [q^e\,x^r\,y^s]\;\bar{J}_{0,N}\,g_{n,n+p,n}
  =(-1)^{r+s}\sum_{w\in\Z}
    \bigl(\sgtwo(r,s)-\sgtwo(U_w,V_w)\bigr)\,
    \mathbf{1}_{e=H(r,s)+N\binom{w+1}{2}}.
\]
\end{corollary}

\begin{proof}
Add Propositions~\ref{prop:s-row} and~\ref{prop:r-row} and use
$\frac{\sg(r)-\sg(U_w)}{2}+\frac{\sg(s)-\sg(V_w)}{2}
=\sgtwo(r,s)-\sgtwo(U_w,V_w)$.
\end{proof}

%----------------------------------------------------------------------
\section{Proof of the Main Identity}\label{sec:proof}
%----------------------------------------------------------------------

\begin{proof}[Proof of Theorem~\ref{thm:main}]
Fix $e,r,s\in\Z$.  Since
$\bar{J}_{0,N}=\sum_{w\in\Z}q^{N\binom{w+1}{2}}$,
the Hecke coefficient is
\[
  [q^e\,x^r\,y^s]\;\bar{J}_{0,N}\,f_{n,n+p,n}
  =(-1)^{r+s}\sum_{w\in\Z}\sgtwo(r,s)\,
    \mathbf{1}_{e=H(r,s)+N\binom{w+1}{2}}.
\]
By Corollary~\ref{cor:g-coeff}:
\[
  [q^e\,x^r\,y^s]\;\bar{J}_{0,N}\,g_{n,n+p,n}
  =(-1)^{r+s}\sum_{w\in\Z}
    \bigl(\sgtwo(r,s)-\sgtwo(U_w,V_w)\bigr)\,
    \mathbf{1}_{e=H(r,s)+N\binom{w+1}{2}}.
\]
Subtracting:
\[
  [q^e\,x^r\,y^s]\bigl(\bar{J}_{0,N}\,f-\bar{J}_{0,N}\,g\bigr)
  =(-1)^{r+s}\sum_{w\in\Z}\sgtwo(U_w,V_w)\,
    \mathbf{1}_{e=H(r,s)+N\binom{w+1}{2}},
\]
which equals $[q^e\,x^r\,y^s]\,\bar\Theta_{n,p}$ by definition.
\end{proof}

%----------------------------------------------------------------------
\section{Examples}\label{sec:examples}
%----------------------------------------------------------------------

\subsection{\texorpdfstring{$f_{2,3,2}$ ($n=2$, $p=1$, $\Delta=5$)}{f232 (n=2, p=1)}}
Here $N=10$, $B=3$, and the Appell sum has $2n=4$ rows.
The main identity gives
$\bar{J}_{0,10}\,f_{2,3,2}=\bar{J}_{0,10}\,g_{2,3,2}+\bar\Theta_{2,1}$.
In the Lean~4 formalization, the right-hand side is packaged as
$\text{Ghat}(1)=g+\theta/\bar{J}$, and the identity
$f_{2,3,2}=\text{Ghat}(1)$ is proved with zero \texttt{sorry}.

\subsection{\texorpdfstring{$f_{1,3,1}$ ($n=1$, $p=2$, $\Delta=8$)}{f131 (n=1, p=2)}}
Here $N=8$ and the Appell sum has $2$ rows.  The cleared theta
correction $\bar\Theta_{1,2}$ is nonzero.  For instance, at
$(r,s)=(1,1)$: $U_0=1$, $V_0=0$, $U_{-1}=0$, $V_{-1}=1$, giving
$[q^3\,xy]\,\bar{J}_{0,8}\,f_{1,3,1}=2$,
$[q^3\,xy]\,\bar{J}_{0,8}\,g_{1,3,1}=0$, and
$[q^3\,xy]\,\bar\Theta_{1,2}=\sgtwo(1,0)+\sgtwo(0,1)=2$.

\subsection{\texorpdfstring{$f_{1,4,1}$ ($n=1$, $p=3$, $\Delta=15$)}{f141 (n=1, p=3)}}
Here $N=15$ and the Appell sum has $2$ rows.
The direct coefficient proof applies without
modification.  The correction $\bar\Theta_{1,3}$ is nonzero:
at $(r,s)=(1,1)$, $[q^4\,xy]\,\bar\Theta_{1,3}=2$, matching the
defect $[q^4\,xy](\bar{J}\cdot f-\bar{J}\cdot g)=2-0=2$.

\subsection{\texorpdfstring{$f_{2,4,2}$ ($n=2$, $p=2$, $\gcd=2$)}{f242 (n=2, p=2, non-coprime)}}
Here $\gcd(n,p)=2$, $N=24$, $B=4$, $\Delta=12$, and the
Appell sum has $2n=4$ rows.  By the $q$-scaling identity
(Section~\ref{sec:discussion}),
$f_{2,4,2}(x,y,q)=f_{1,2,1}(x,y,q^2)$, so the Hecke sum
admits both the $(2,2)$-decomposition of
Theorem~\ref{thm:main} and the scaled primitive
$(1,1)$-decomposition.  The coefficient proof applies to
the $(2,2)$-decomposition directly: no step uses
$\gcd(n,p)=1$.

%----------------------------------------------------------------------
\section{Formal Verification}\label{sec:lean}
%----------------------------------------------------------------------

Table~\ref{tab:lean} summarizes the Lean~4 formalization status.

\begin{table}[ht]
\centering
\begin{tabular}{lll}
\hline
\textbf{Claim} & \textbf{Status} & \textbf{File / theorem} \\
\hline
$f_{2,3,2}$ (spiral param.) & fully verified & \texttt{Chapter10\_M3\_AppellCoeffNF} \\
Row-exponent identity & verified $\forall n,p$ & \texttt{Chapter10\_RootPairGeneral} \\
Branch-sign lemma & verified $\forall n,p$ & \texttt{Chapter10\_RootPairGeneral} \\
Parity criterion & verified $\forall n,p$ & \texttt{Chapter10\_RootPairGeneral} \\
Props.~\ref{prop:s-row},\ref{prop:r-row} (general) & fully verified & \texttt{Chapter10\_GeneralHahnCoeff} \\
Main identity, all $n,p\ge1$ & fully verified & \texttt{general\_boundary\_cleared\_hm\_coefficient} \\
Kronecker coefficient proof & fully verified & \texttt{Chapter10\_KroneckerCoeff} \\
HM residue comparison & fully verified & \texttt{Chapter10\_HMComparison} \\
\hline
\end{tabular}
\caption{Lean~4 formalization coverage.  All verified items have
zero \texttt{sorry} and no axioms beyond \texttt{propext},
\texttt{Classical.choice}, \texttt{Quot.sound}.}
\label{tab:lean}
\end{table}

The $f_{2,3,2}$ verification (commit \texttt{5ead660})
proves the identity
\texttt{Fhm\_eq\_Ghat\_via\_uvOmega} via a wall-coordinate
specialization: parity split of the 4~Appell rows, finsum
reindex via coordinate equivalences, wall-defect decomposition,
and root-pair cancellation.  The Lean object \texttt{Ghat}
is defined as a four-row Appell sum specialized to HM's
monomial-spiral parameter (cf.~\cite{jHiMo14}, Theorem~0.6).
This older regression uses a different specialization interface and is
logically independent of the general coefficient theorem formalized here.

The general algebraic core is formalized in
\nolinkurl{Chapter10_RootPairGeneral.lean};
\nolinkurl{Chapter10_GeneralHahnCoeff.lean} adds the explicit finite residue
rows, their monomial-forcing collapse, local finiteness, both row formulas,
and the main coefficient assembly.  The appendix is formalized independently:
\nolinkurl{Chapter10_KroneckerCoeff.lean} proves the PF quasi-period by an
inverse-coefficient telescope, and \nolinkurl{Chapter10_HMComparison.lean}
retains both residue variables and all three bilateral indices in a constrained
source type, proves its equivalence with the single shell index $w$, and checks
the branch-sign and exponent formulas (A.10)--(A.17).  The final module is
\nolinkurl{Chapter10_HahnHeckeAppell}; its main theorem is
\nolinkurl{general_boundary_cleared_hm_coefficient}.

The self-contained fixed registry source, its exact proof closure, and the
paper snapshot are archived together at
\url{https://github.com/xiangyazi24/palomar-hahn-hecke-appell}.
The relevant Lean files are in the \texttt{Pending/} directory and the final
audit target is
\nolinkurl{QseriesFormalization.Pending.Chapter10_HahnHeckeAppellAudit}.
A Sage script verifying both the row-exponent identity
(Proposition~\ref{prop:exponent}) and the appendix exponent
identity~(A.15) is included as \texttt{check\_exponents.sage}.

%----------------------------------------------------------------------
\section{Discussion}\label{sec:discussion}
%----------------------------------------------------------------------

\subsection*{Comparison with existing proofs}

The Hickerson--Mortenson proof~\cite{jHiMo14} uses analytic
uniqueness; Mortenson~\cite{jMo12} uses ${}_1\psi_1$.  Our proof
replaces both with Hahn-series coefficient algebra.  The $\sgtwo$
structure arises from the Hahn geometric inverse expansion direction,
providing a formal-algebraic interpretation of the Appell--Lerch
sign-wall coefficients.

The row-exponent identity (Proposition~\ref{prop:exponent}) is the
algebraic heart: it collapses the multi-index Appell row exponent
into the Hecke exponent $H(r,s)$ plus a theta shell $N\binom{w+1}{2}$.
This identity is verifiable by \texttt{ring} and is the key step
that makes the coefficient comparison self-contained.

We note the relationship with Bailey-pair approaches to Hecke-type
double sums~\cite{jAn01,jLO13}.  Those methods are also
formal-series arguments and produce identities for one-variable
specializations; however, they do not establish the two-variable
$(x,y)$-generic Appell--Lerch \emph{decomposition}, where the
bilateral Appell parameter is precisely where analytic or Hahn-algebraic
machinery enters.

\subsection*{Non-coprime cases}

For $\gcd(n,p)=d>1$, one has
$f_{dn_0,d(n_0+p_0),dn_0}(x,y;q)=f_{n_0,n_0+p_0,n_0}(x,y;q^d)$,
so the Hecke sum admits both the direct $(n,p)$-decomposition
(Theorem~\ref{thm:main}) and the scaled primitive
$(n_0,p_0)$-decomposition.  The two decompositions have different
numbers of Appell rows and different base parameters; their
coexistence is a consistency check, not a derivation of one
from the other.

\subsection*{Connection to indefinite theta functions and wall-crossing
  (formal analogy, not used in the proof)}

The $\sgtwo$ kernel appearing in $\bar\Theta_{n,p}$ is of the
same type as the sign-wall functions arising in Zwegers'
indefinite theta series~\cite{oZw02} and in the wall-crossing
analysis of Dabholkar--Murthy--Zagier~\cite{jDMZ14}.
The connection is structural: the Hecke exponent $H(r,s)$
is an indefinite quadratic form of type~$(1,1)$---the same
structure that Hecke~\cite{oHe25} studied in his original
construction of modular forms from indefinite quadratic
forms---and the chosen Hahn lexicographic order is the
formal expansion chamber analogous to the analytic chamber
in which the corresponding monomials satisfy $|q|<1$.

From a polyhedral viewpoint, the signed Hahn geometric inverse
(Lemma~\ref{lem:hahn}) is a one-dimen\-sional instance of the
polarized-cone expansion of
Varchenko~\cite{jVa87} and Brion~\cite{jBr88}: expanding
$1/(1-M)$ in one Hahn chamber selects one ray of a
one-dimensional cone, and the branch sign records the
orientation.  The theta correction $\bar\Theta_{n,p}$ is then
the boundary term arising from the polarization.  The use of
ordered formal series rings to select expansion branches is
also developed by Xin~\cite{jXi05} (Malcev--Neumann series
and formal residues) and Aparicio
Monforte--Kauers~\cite{jAMK13} (formal Laurent series in
several variables).

What distinguishes the formal
setting is that the $\sgtwo$ kernel remains exact---no Vign\'eras
error function~\cite{oVi77} or modular completion is needed,
because the Hahn ring has no $\bar\tau$ variable to smooth.
This connection is included as context; the paper does not prove
a modular completion theorem.

\subsection*{Coefficient data for quantum-modular settings
  (possible application, not proved here)}

The paper formulas are integer-valued coefficient identities.  The current
Lean wrapper checks them over $\Q$ and does not separately formalize
integrality or arbitrary-ring base change.  The resulting coefficient data
is available as input to quantum-modular and radial-limit
analysis~\cite{jZag10,jFOR13}, where the analytic and
${}_1\psi_1$ proofs do not directly supply such data.

\subsection*{Extensions}

The approach may extend to other families of indefinite theta
identities.  For $f_{a,b,c}$ with $a\neq c$ (HM Theorem~0.4),
the row-exponent identity generalizes, but the branch-sign lemma
requires a case analysis whose boundary cases encode the
additional correction structure.  The key requirement is a
row-exponent collapse analogous to Proposition~\ref{prop:exponent};
when such a collapse exists, the branch-sign mechanism produces the
correction term automatically.

\bibliographystyle{plain}
\bibliography{master}

%----------------------------------------------------------------------
\appendix
\section{Coefficient comparison with the Hickerson--Mortenson theta term}
\label{app:HM-coeff-comparison}
%----------------------------------------------------------------------

We keep the conventions of the paper:
\[
  C_2(z)=\binom z2=\frac{z(z-1)}2,\qquad
  B=n+p,\qquad
  \Delta=p(2n+p),\qquad
  N=n\Delta,
\]
\[
  A_0=\frac{np(2n+p+1)}2,\qquad
  H(R,S)=nC_2(R)+BRS+nC_2(S),\qquad
  \sgtwo(u,v)=\frac{\sg(u)+\sg(v)}2,
\]
and
\[
  j(z;q)=\sum_{m\in\mathbb Z}(-1)^m q^{C_2(m)}z^m,\qquad
  \sg(0)=1.
\]
Let
\[
  \alpha=\frac{n-1}{2},\qquad \beta=\frac{n+1}{2},\qquad
  \varepsilon=\{\alpha\}\in\{0,1/2\}.
\]
For \(r^*,s^*\in\{0,\ldots,p-1\}\), put
\[
  \widetilde r=r^*+\varepsilon,\qquad
  \widetilde s=s^*+\varepsilon,\qquad
  \rho=\widetilde r-\alpha,\qquad
  \sigma=\widetilde s+\beta.
\]
Then \(\rho,\sigma\in\mathbb Z\), and
\(\widetilde r,\widetilde s\in \varepsilon+\{0,\ldots,p-1\}\subset[0,p)\).
Also
\[
  \widetilde s-\widetilde r=\sigma-\rho-n.
\]
The Hickerson--Mortenson theta correction, specialized to
\((z_0,z_1)=(-1,-1)\), can be rewritten as
\[
\begin{aligned}
\theta^{\rm HM}_{n,p}
=&\sum_{r^*=0}^{p-1}\sum_{s^*=0}^{p-1}
 q^{H(\rho,\sigma)}(-x)^\rho(-y)^\sigma
 j\!\left(-q^{np(\sigma-\rho-n)}\frac{x^n}{y^n};q^{np^2}\right) \\
&\qquad\qquad\times
 \frac{J_Q^3\,j(d_1d_2;Q)}{j(d_1;Q)j(d_2;Q)},
\end{aligned}
\tag{A.1}
\]
where
\[
  Q=q^{p\Delta}=q^{p^2(2n+p)},
\]
\[
  d_1=q^{\Delta\rho+A_0-np}\frac{(-y)^B}{(-x)^n},\qquad
  d_2=q^{\Delta\sigma-A_0}\frac{(-x)^B}{(-y)^n}.
\]
Indeed,
\[
  d_1d_2=q^{\Delta(\rho+\sigma)-np}(-x)^p(-y)^p,
\]
which is the second numerator theta argument in HM's formula, since
\[
  \Delta(\widetilde r+\widetilde s)+pB
  =\Delta(\rho+\sigma-1)+pB
  =\Delta(\rho+\sigma)-np.
\]

\begin{lemma}[Kronecker partial fraction]\label{lem:kronecker}
With $j(z;Q)=\sum_m(-1)^mQ^{\binom{m}{2}}z^m$,
\[
  \frac{J_Q^3\,j(ab;Q)}{j(a;Q)\,j(b;Q)}
  =\sum_{\eta\in\mathbb Z}\frac{b^\eta}{1-aQ^\eta}.
\tag{A.2}
\]
\end{lemma}

\begin{proof}
We give the coefficientwise proof used in the Lean development.  Put
\[
  \operatorname{PF}(a)=
  \sum_{k\in\mathbb Z}
    \frac{(-1)^kQ^{\binom{k+1}{2}}}{1-aQ^k},
\]
with every denominator expanded in its signed $Q$-first Hahn branch.
The coefficientwise partial-fraction theorem
\[
  j(a;Q)\operatorname{PF}(a)=J_Q^3
\]
follows by a geometric telescope: all nonconstant $a$-coefficients cancel,
while the constant coefficient is Jacobi's cube identity.

Let $K(a,b)=\sum_{\eta\in\mathbb Z}b^\eta/(1-aQ^\eta)$.
The coefficient of $b^B$ in $j(b;Q)K(a,b)$ is
\[
  \sum_{s\in\mathbb Z}
    \frac{(-1)^sQ^{\binom{s}{2}}}{1-aQ^{B-s}}
  =\operatorname{PF}(aQ^B),
\]
where the second equality uses $k=-s$ and
$\binom{s}{2}=\binom{k+1}{2}$.  The inverse-coefficient telescope also gives
\[
  \operatorname{PF}(aQ^B)
  =(-1)^BQ^{\binom{B}{2}}a^B\operatorname{PF}(a)
  \qquad(B\in\mathbb Z).
\]
In the one-step relation, the only boundary term is a coefficient of
$j(1;Q)$, whose two roots $k$ and $-k-1$ cancel.  The last display is exactly
the coefficient of $b^B$ in $j(ab;Q)\operatorname{PF}(a)$.  Hence
\[
  j(b;Q)K(a,b)=j(ab;Q)\operatorname{PF}(a)
\]
coefficientwise.  Multiplication by $j(a;Q)$ and the partial-fraction theorem
above gives
\[
  j(a;Q)j(b;Q)K(a,b)=J_Q^3j(ab;Q);
\]
cancellation in the
corresponding Hahn fraction field proves (A.2).  This is the classical
Kronecker identity (cf.~Weil~\cite{oWe76}, Chapter~IV,
and~\cite{jHiMo14}, Proposition~5.7 and Lemma~5.10), but no residue or
Liouville argument is needed here.
\end{proof}
Applying Lemma~\ref{lem:kronecker} with $a=d_1$, $b=d_2$ gives
\[
  \frac{J_Q^3j(d_1d_2;Q)}{j(d_1;Q)j(d_2;Q)}
  =\sum_{\eta\in\mathbb Z}\frac{d_2^\eta}{1-Q^\eta d_1}.
\tag{A.3}
\]
Next expand the first numerator theta factor by the Jacobi triple product:
\[
\begin{aligned}
 j\!\left(-q^{np(\sigma-\rho-n)}\frac{x^n}{y^n};q^{np^2}\right)
 &=\sum_{\iota\in\mathbb Z}
 (-1)^\iota q^{np^2 C_2(\iota)}
 \left(-q^{np(\sigma-\rho-n)}\frac{x^n}{y^n}\right)^\iota \\
 &=\sum_{\iota\in\mathbb Z}
 q^{np^2 C_2(\iota)+np(\sigma-\rho-n)\iota}
 (-x)^{n\iota}(-y)^{-n\iota}.
\end{aligned}
\tag{A.4}
\]
Here the signs cancel because \((-1)^\iota(-1)^\iota=1\).

For the remaining denominator, set
\[
  M_{\rho,\eta}=Q^\eta d_1
  =q^{p\Delta\eta+\Delta\rho+A_0-np}\frac{(-y)^B}{(-x)^n}.
\]
Its leading \(q\)-exponent is
\[
  \Gamma_{\rho,\eta}=p\Delta\eta+\Delta\rho+A_0-np.
\]
\emph{Formal--analytic bridge.}
On the unit torus $|x|=|y|=1$ with $|q|<1$,
$|M_{\rho,\eta}|=|q|^{\Gamma_{\rho,\eta}}$, so
$|M_{\rho,\eta}|<1$ iff $\Gamma_{\rho,\eta}>0$.
Hence the Hahn branch of $(1-M_{\rho,\eta})^{-1}$
coincides with the convergent geometric series.  Thus the coefficientwise
proof of Lemma~\ref{lem:kronecker} specializes to the analytic quotient on
this torus, with the same Hahn branch convention.
The signed Hahn geometric expansion gives
\[
  (1-M_{\rho,\eta})^{-1}
  =\sum_{\ell\in\mathbb Z} c_{\rho,\eta}(\ell)M_{\rho,\eta}^{\ell},\qquad
  c_{\rho,\eta}(\ell)=\frac{\sg(\ell)+\sg(\Gamma_{\rho,\eta})}{2}.
\tag{A.5}
\]

Combining (A.1)--(A.5), a term indexed by
\((\rho,\sigma,\iota,\eta,\ell)\) has monomial part
\[
  (-x)^{\rho+n\iota+B\eta-n\ell}
  (-y)^{\sigma-n\iota-n\eta+B\ell}.
\]
Thus the coefficient of \((-x)^R(-y)^S\) is constrained by
\[
  R=\rho+n\iota+B\eta-n\ell,\qquad
  S=\sigma-n\iota-n\eta+B\ell.
\tag{A.6}
\]
Define
\[
  w=\ell-\iota-\eta.
\tag{A.7}
\]
Then (A.6) implies
\[
  U_w:=R+nw+\alpha=\rho+p\eta+\alpha=\widetilde r+p\eta,
\tag{A.8}
\]
\[
  V_w:=S-nw-\beta=\sigma-\beta+p\ell=\widetilde s+p\ell.
\tag{A.9}
\]
Conversely, for fixed \(R,S,w\), equations (A.8)--(A.9) determine the
contributing residue pair and summation indices uniquely.  Indeed, since
\(\widetilde r\in\varepsilon+\{0,\ldots,p-1\}\), there is a unique
\(\widetilde r\) in that residue interval and a unique \(\eta\in\mathbb Z\)
such that
\[
  R+nw+\alpha=\widetilde r+p\eta.
\]
Similarly, there is a unique \(\widetilde s\in\varepsilon+\{0,\ldots,p-1\}\)
and a unique \(\ell\in\mathbb Z\) such that
\[
  S-nw-\beta=\widetilde s+p\ell.
\]
Then \(\rho=\widetilde r-\alpha\), \(\sigma=\widetilde s+\beta\), and
(A.7) determines the unique remaining index
\[
  \iota=\ell-\eta-w.
\]
Thus coefficient extraction over the HM residue variables and the bilateral
indices is equivalent to summing over \(w\in\mathbb Z\).

We now identify the branch coefficient.  Since
\(\widetilde s\in[0,p)\), equation (A.9) gives
\[
  \sg(\ell)=\sg(\widetilde s+p\ell)=\sg(V_w).
\tag{A.10}
\]
Indeed this is immediate for \(\ell\ge0\), while for \(\ell\le -1\) one has
\(\widetilde s+p\ell<0\).  Likewise, using (A.8),
\[
\begin{aligned}
  \Gamma_{\rho,\eta}
  &=p\Delta\eta+\Delta\rho+A_0-np \\
  &=\Delta(\rho+p\eta)+A_0-np \\
  &=\Delta(U_w-\alpha)+A_0-np \\
  &=\Delta U_w+\frac{pB}{2}.
\end{aligned}
\tag{A.11}
\]
If \(U_w\ge0\), then \(\Gamma_{\rho,\eta}>0\).  If \(U_w<0\), then
\(U_w=\widetilde r+p\eta\) with \(\eta\le -1\), hence
\(U_w\le \widetilde r-p\le -1/2\), and therefore
\[
  \Gamma_{\rho,\eta}\le -\frac{\Delta}{2}+\frac{pB}{2}
  =-\frac{np}{2}<0.
\]
Thus
\[
  \sg(\Gamma_{\rho,\eta})=\sg(U_w).
\tag{A.12}
\]
Combining (A.5), (A.10), and (A.12), the branch coefficient becomes
\[
  c_{\rho,\eta}(\ell)
  =\frac{\sg(V_w)+\sg(U_w)}2
  =\sgtwo(U_w,V_w).
\tag{A.13}
\]

It remains to compute the \(q\)-exponent.  The total exponent of the term
indexed by \((\rho,\sigma,\iota,\eta,\ell)\) is
\[
\begin{aligned}
\mathcal E
=&H(\rho,\sigma)
 +np^2C_2(\iota)+np(\sigma-\rho-n)\iota \\
&+(\Delta\sigma-A_0)\eta
 +(p\Delta\eta+\Delta\rho+A_0-np)\ell.
\end{aligned}
\tag{A.14}
\]
Using (A.6) and (A.7), a direct expansion gives
\[
  \mathcal E=H(R,S)+N C_2(w+1).
\tag{A.15}
\]
The following Sage snippet is included as an audit aid; the
identity is an elementary polynomial expansion:
\begin{verbatim}
var('rho sigma iota eta ell n p')
C2 = lambda z: z*(z-1)/2
B = n + p
Delta = p*(2*n + p)
N = n*Delta
A0 = n*p*(2*n + p + 1)/2
R = rho + n*iota + B*eta - n*ell
S = sigma - n*iota - n*eta + B*ell
w = ell - iota - eta
H = lambda X,Y: n*C2(X) + B*X*Y + n*C2(Y)
E = H(rho,sigma) + n*p^2*C2(iota) + n*p*(sigma-rho-n)*iota \
    + (Delta*sigma - A0)*eta \
    + (p*Delta*eta + Delta*rho + A0 - n*p)*ell
expand(E - H(R,S) - N*C2(w+1))
# output: 0
\end{verbatim}

Therefore, for fixed \(e,R,S\), equations (A.13)--(A.15) give
\[
\begin{aligned}
\relax[q^e(-x)^R(-y)^S]\,\theta^{\rm HM}_{n,p}
=\sum_{w\in\mathbb Z}
 \sgtwo\!\left(R+nw+\frac{n-1}{2},
               S-nw-\frac{n+1}{2}\right)
 \mathbf 1_{\,e=H(R,S)+N C_2(w+1)}.
\end{aligned}
\tag{A.16}
\]
Equivalently, since \((-x)^R(-y)^S=(-1)^{R+S}x^Ry^S\),
\[
\begin{aligned}
\relax[q^e x^R y^S]\,\theta^{\rm HM}_{n,p}
=&(-1)^{R+S}
\sum_{w\in\mathbb Z}
 \sgtwo\!\left(R+nw+\frac{n-1}{2},
               S-nw-\frac{n+1}{2}\right)  \\
&\qquad\qquad\times
 \mathbf 1_{\,e=H(R,S)+N C_2(w+1)}.
\end{aligned}
\tag{A.17}
\]
The sum is locally finite: for fixed \(e,R,S\), the equation
\(e=H(R,S)+N C_2(w+1)\) has at most two integer solutions \(w\).

But the coefficient of \(q^e x^R y^S\) in
\[
\bar\Theta_{n,p}
=\sum_{R,S,w\in\mathbb Z}(-x)^R(-y)^S
 \sgtwo\!\left(R+nw+\frac{n-1}{2},
               S-nw-\frac{n+1}{2}\right)
 q^{H(R,S)+N C_2(w+1)}
\]
is exactly the right-hand side of (A.17).  Hence
\[
  \theta^{\rm HM}_{n,p}(x,y,q)=\bar\Theta_{n,p}(x,y,q)
\]
coefficientwise.


\end{document}
